\section{Related work}\label{sec:related}

Much of the skeleton of this setting is classical. Markov logic networks tie rules by role across all their groundings, fit them by pseudolikelihood, and solve them by
local search that stops when every hard clause holds\cite{richardson2006markov}. Their authors also noted that pseudolikelihood can do badly when long chains of
inference are needed\cite{lowd2007efficient}. One conditional model per relation, with counts over related objects as
features, is a relational dependency network\cite{neville2007relational}, and counting the third objects that pass a test
on both legs is one of the complex aggregates of relational decision trees\cite{vens2006refining}. Asking that each
solved example be a model of the learned rules is learning from interpretations\cite{deraedt1994first}. Programs with
recursion are learned exactly from examples by inductive logic programming: ILASP learns answer set programs whose
examples must be answer sets\cite{law2014inductive}; Popper searches for recursive programs\cite{cropper2021learning};
differentiable forward chaining learns rules by unrolling inference\cite{evans2018learning}. ``Derive what is forced,
then guess and check'' is how answer set solvers work, starting from the well-founded model\cite{vangelder1991well}.
Constraint acquisition learns count constraints from solutions, and COUNT-CP transfers them from $4\times4$ to $9\times9$
Sudoku\cite{kumar2022learning}. Every symbolic learner here searches the space of rules; we fit a fixed family of
conditions in one step.

On the neural side, CRAM is the immediate starting point\cite{hoover2026classical}, and energies for reasoning are an
active line\cite{du2022learning,chen2026energy}. Looped networks trained on small mazes and run on far larger ones go
back to Deep Thinking\cite{bansal2022end}. Recurrent and message-passing reasoners now transfer across sizes on
constraint problems\cite{toenshoff2023one,salvatori2026tiny}, and lattice deduction transformers, which project a
recurrent state onto a lattice at every step, solve all of Sudoku-Extreme\cite{davis2026lattice}. In algorithmic reasoning, hints supervise a network
on an algorithm's trace\cite{velickovic2022clrs}, and a no-hint mode trains on answers alone, through the
loop\cite{rodionov2023neural}. Discrete neural algorithmic reasoning proves its learned executions correct on any input
but trains on hints\cite{rodionov2025discrete}. Dual algorithmic reasoning adds an optimality certificate as extra
supervision\cite{numeroso2023dual}, and a sparsity regulariser lets a network trained on a few shortest-path instances
provably approximate Bellman--Ford at every size\cite{nerem2025graph}. Proof labelling schemes and locally checkable
proofs study, in distributed computing, which global properties have labels that a local verifier can
check\cite{korman2010proof,korman2007distributed,goos2016locally}.

Against this background, our contributions are four. The first is the cold pair read with comparisons of counts,
fitted at the answer in one step with no search over rules. The second is the answer completed from itself, by
closures and counts, as the training target. It links the hints of algorithmic reasoning to the labels of proof
labelling schemes: a completion is a certificate computed from the answer instead of from the trace. The third is the
level condition, kept as a safeguard, which turns a check fitted at a supported model into a procedure that computes it
wherever the fit alone would drop it. The
fourth is proofs from the fitted table at every size. Our searches, whose scope Methods describes, found none of these
four in a learned system.
